\chapter{Sequences, recursion, and discrete asymptotics}

\begin{orient}
\textbf{What this chapter is for.} A recursively defined sequence $x_{n+1}=f(x_n)$ poses two entirely separate questions, and almost every wrong answer in this area comes from answering only the second. The first question is whether the sequence converges at all; the second is what it converges to. Existence is settled by an inequality argument (monotone and bounded, or a contraction estimate); the value is then settled by solving $L=f(L)$ and discarding the roots the sequence can never reach. The second half of the chapter turns to the asymptotic side: how to extract the growth rate of a sequence you cannot solve, using Stolz-Cesaro, averaging, $n$th-root and ratio limits, and Stirling's approximation. By the end you should be able to look at a recursion and name, within a few seconds, which existence argument it wants, and to read the size of $a_n$ off its shape when no closed form is available.
\end{orient}

\section{Two questions, not one}

Here is the error this chapter exists to kill. A student is given $x_1=1$ and $x_{n+1}=\sqrt{2+x_n}$ and asked for $\lim x_n$. They write: let $L$ be the limit, then $L=\sqrt{2+L}$, so $L^2-L-2=0$, so $L=2$. They box the answer and move on. The algebra is correct and the answer is correct, and the argument is worth almost nothing, because it never established that a limit exists. The identical manipulation applied to $x_1=1$, $x_{n+1}=3x_n-2$ yields $L=3L-2$, hence $L=1$; but that sequence is constantly $1$ only if it starts at $1$, and from $x_1=2$ it runs away to $+\infty$ while the fixed-point equation sits there unchanged, still offering $L=1$. The equation $L=f(L)$ is a \emph{consequence} of convergence, not a cause of it.

So the discipline is fixed and it does not vary from problem to problem. Step one: prove that $\lim x_n$ exists, by an argument that never mentions its value. Step two: use continuity of $f$ to pass to the limit in $x_{n+1}=f(x_n)$, obtaining $L=f(L)$, and then select among its roots using what step one told you about the range of the sequence. Step one is where the mathematics is. Step two is bookkeeping.

There is a practical corollary about how to write the answer. The sentence ``let $L=\lim x_n$'' is a claim, not a definition, and it should not appear on the page until it has been earned. Write the two inductions first, then the sentence, then the algebra. Examiners award the marks in that order because the mathematics is in that order, and a script that opens with ``let $L$ be the limit'' has spent its first line asserting the thing it was asked to prove.

\begin{keynote}[EXISTENCE BEFORE VALUE]
$L=f(L)$ is what a limit must satisfy if it exists. It is silent on whether it exists. A solution that solves the fixed-point equation and stops is not a partial solution; it is a solution to a different problem.
\end{keynote}

Why is step one enough? Because of completeness. A bounded increasing sequence has a supremum $S$, and for any $\varepsilon>0$ some term exceeds $S-\varepsilon$; monotonicity then keeps every later term in $(S-\varepsilon,S]$, so $x_n\to S$. That is the whole proof, and it is worth carrying in your head, because it explains why the theorem gives you no information about the value: the supremum is produced by an axiom, not by a formula. The fixed-point equation is the only route to the number, and it is available only once the supremum is known to be a limit.

The bound you need is never a mystery either. Solve $L=f(L)$ first, informally, as reconnaissance; the roots tell you where the sequence is trying to go, and the relevant root is the bound to attempt in the induction. Doing the algebra first is fine. Presenting it first, as the argument, is not.

There are only two existence arguments in wide use, and they split by the shape of $f$. If $f$ is increasing on the interval where the sequence lives, the sequence is automatically monotone, and you need only a bound: that is the monotone convergence theorem. If $f$ is not increasing, or the bound is awkward, you instead show that $f$ shrinks distances near the fixed point, which is the contraction argument. Everything else in the first half of this chapter is a refinement of one of those two.

\tech{Monotone convergence: proving that the limit exists}{easy}
\cue A recursion $x_{n+1}=f(x_n)$ with no closed form, where $f$ is increasing on the range of the sequence, and the problem asks for $\lim x_n$.
\method Prove two statements by induction, separately: (i) the sequence is bounded (guess the bound from the fixed point), and (ii) it is monotone. A bounded monotone real sequence converges; this is the completeness of $\R$ in its most usable form. Only after both inductions are complete may you write $L=f(L)$.

For monotonicity, the cleanest route when $f$ is increasing is to check the first step only. If $x_2>x_1$, then applying the increasing $f$ to both sides gives $x_3>x_2$, and induction carries it forever. The sign of $x_2-x_1$ decides the direction of the whole sequence.

\begin{worked}
Let $x_1=1$ and $x_{n+1}=\sqrt{2+x_n}$.

\emph{Bounded above by $2$.} True for $x_1$. If $x_n<2$ then $x_{n+1}=\sqrt{2+x_n}<\sqrt{4}=2$.

\emph{Increasing.} Since all terms are positive, compare squares:
\[x_{n+1}^2-x_n^2=2+x_n-x_n^2=(2-x_n)(1+x_n)>0\]
because $0<x_n<2$. So $x_{n+1}>x_n$.

Bounded and increasing, so $L=\lim x_n$ exists and lies in $[1,2]$. Now $f(t)=\sqrt{2+t}$ is continuous, so $L=\sqrt{2+L}$, giving $L^2-L-2=(L-2)(L+1)=0$. The root $L=-1$ is impossible since $L\geq x_1=1$. Hence $L=2$.

\emph{The same recursion from above.} Start instead at $x_1=3$. Now $x_{n+1}^2-x_n^2=(2-x_n)(1+x_n)<0$, so the sequence decreases, and $x_n>2$ is preserved: $x_{n+1}=\sqrt{2+x_n}>\sqrt4=2$. Decreasing and bounded below, so it converges, and the limit still satisfies $L=\sqrt{2+L}$, giving $L=2$ again. The terms are $2.2361$, $2.0582$, $2.0145$, $2.0036$, $2.0009$. The lesson: $f$ increasing means the sign of $x_2-x_1$ sets the direction, and the fixed point is approached from whichever side you started on.
\end{worked}

\begin{trap}
Solving $L=f(L)$ with no existence argument. The fixed-point equation for $x_{n+1}=3x_n-2$ has the root $L=1$ for every starting value, including the ones that diverge. A second, quieter failure: proving boundedness and forgetting monotonicity, which leaves open the oscillating case.
\end{trap}

\begin{seealsob}Contraction and the derivative criterion; subsequence arguments; solving the fixed-point equation and selecting the admissible root.\end{seealsob}

Notice how little of that argument was computation. Two inductions, one continuity remark, one root rejected on sign: the entire content is inequalities. This is typical, and it is why the technique is tiered easy despite being the load-bearing result of the chapter. What makes it fail in practice is not difficulty but omission.

Once existence is in hand, the remaining work is to choose correctly among the roots of $L=f(L)$. That sounds trivial, and for a quadratic with one positive and one negative root it is. It stops being trivial when both roots are admissible-looking, which is exactly the situation the next technique is built for. Note also that a great many problems are \emph{posed} in the self-similar form: nested radicals, periodic continued fractions and towers are all presented as static objects rather than as sequences, and the first move is always to name the sequence of truncations whose limit the object is defined to be. Until you have done that there is nothing to prove convergence of.

\tech[Solving the fixed point equation]{Solve $L=f(L)$, then select the admissible root}{medium}
\cue A recursion whose tail reproduces the whole expression: a periodic continued fraction, an infinite nested radical, an iterated exponential tower, a self-similar product. The defining feature is that the object contains a copy of itself.
\method Name the value $L$, write the self-similarity as $L=f(L)$, and solve. Then discard every root that the sequence's own range forbids, using the bound you proved in step one. Sign, monotonicity and the interval $[\inf x_n,\sup x_n]$ are the three selectors.

\begin{worked}
Evaluate the tower $t^{t^{t^{\cdots}}}$ at $t=\sqrt2$, meaning $\lim x_n$ for $x_1=\sqrt2$, $x_{n+1}=t^{x_n}$.

\emph{Existence.} $f(u)=(\sqrt2)^{u}$ is increasing. If $x_n\leq2$ then $x_{n+1}\leq(\sqrt2)^2=2$, so the sequence is bounded above by $2$; and $x_2=(\sqrt2)^{\sqrt2}\approx1.6325>\sqrt2=x_1$, so it increases. The limit $L$ exists and $L\in[\sqrt2,2]$.

\emph{Value.} $L=(\sqrt2)^{L}$, that is $L^{1/L}=\sqrt2$. This equation has \emph{two} positive roots: $L=2$, since $2^{1/2}=\sqrt2$, and $L=4$, since $4^{1/4}=\sqrt2$. The bound from step one kills $L=4$ outright. Hence $L=2$; numerically the iteration from $x_1=\sqrt2$ reaches $2.0000000000000004$ in double precision.

\emph{A periodic continued fraction.} Evaluate $x=1+\cfrac{1}{2+\cfrac{1}{2+\cdots}}$, that is the limit of the convergents of $[1;2,2,2,\ldots]$. Writing $y$ for the repeating tail $2+1/(2+\cdots)$ gives $y=2+1/y$, so $y^2-2y-1=0$ and $y=1+\sqrt2$ (the negative root is impossible since every convergent exceeds $2$). Then $x=1+1/y=1+(\sqrt2-1)=\sqrt2$. The convergents are alternately above and below, so each consecutive pair brackets the answer, and the iteration confirms $1.41421356237$.
\end{worked}

\begin{trap}
Reporting both roots, or the wrong one. The tower at $t=\sqrt2$ is the standard demonstration: $4$ satisfies the equation and is not the limit, because the sequence never gets past $2$. The fixed-point equation is necessary and not sufficient, and it is the sequence's range, established independently, that does the selecting.
\end{trap}

\begin{seealsob}Monotone convergence; contraction and the derivative criterion; subsequence arguments.\end{seealsob}

\subsection{Recursions that are secretly explicit}

Before reaching for any existence theorem, spend ten seconds asking whether the recursion can be solved outright, because a closed form makes every question in this chapter trivial. Three substitutions account for most of the solvable cases.

\emph{Reciprocals linearise Moebius recursions.} If $x_{n+1}=\dfrac{x_n}{1+x_n}$ then $\dfrac{1}{x_{n+1}}=\dfrac{1}{x_n}+1$, so $1/x_n$ is an arithmetic progression: from $x_1=1$ we get $1/x_n=n$ and $x_n=1/n$ exactly. No induction, no fixed point, no rate estimate; the answer is $0$ and the rate is $1/n$. More generally any $x_{n+1}=(ax_n+b)/(cx_n+d)$ becomes linear in $1/(x_n-\lambda)$ where $\lambda$ is a fixed point.

\emph{Logarithms linearise power recursions.} If $x_{n+1}=x_n^2$ then $\ln x_{n+1}=2\ln x_n$, so $x_n=x_1^{2^{n-1}}$: the sequence tends to $0$ for $\abs{x_1}<1$, to $1$ at $x_1=1$, and to $\infty$ beyond. The fixed-point equation $L=L^2$ offers $0$ and $1$ and cannot distinguish these three regimes on its own.

\emph{Trigonometric substitution linearises Chebyshev recursions.} If $x_{n+1}=2x_n^2-1$ and $\abs{x_1}\leq1$, write $x_1=\cos\theta$; the double-angle formula gives $x_n=\cos(2^{n-1}\theta)$ exactly, as the values $0.7648$, $0.1700$, $-0.9422$, $0.7756$ from $\theta=0.7$ confirm. The sequence therefore has no limit for generic $\theta$, and the reason is visible: the angle doubles forever. The fixed-point equation $L=2L^2-1$ has the perfectly good roots $1$ and $-\tfrac12$, neither of which the sequence approaches.

That last example is the sharpest warning in the chapter. A recursion can be completely understood, have two real fixed points, and converge to neither.

\section{Rates, and the local picture at a fixed point}

The self-similarity move has one further use worth stating, because it is where the technique earns its hard tier: it applies to objects that are not obviously sequences at all. A doubly infinite product, an integral defined by a functional equation, a probability defined by conditioning on the first step, all reduce to $L=f(L)$ plus a range argument. The range argument is the part that gets skipped, and it is the part that is doing the work.

Knowing that $x_n\to L$ is often not the end of the question. Numerical work needs to know how many steps buy a decimal digit; qualitative work needs to know whether the approach is one-sided or oscillating, since a one-sided approach means every term is a bound on the answer. Both are read off the derivative of $f$ at the fixed point, and both can be seen before any calculation by drawing the cobweb.

\tech{Reading a recursion off the cobweb}{easy}
\cue You want the qualitative behaviour of $x_{n+1}=f(x_n)$ (monotone, alternating, escaping) before committing to an inductive proof, or the induction is not obviously going to close.
\method Draw $y=f(x)$ and $y=x$ on the same axes. Starting at $x_1$ on the horizontal axis, go vertically to the curve (this is $x_2$), horizontally to the diagonal (this transfers $x_2$ to the horizontal axis), and repeat. The path staircases or spirals; the fixed points are the crossings, and the slope of $f$ at a crossing dictates the local picture. For $0<f'(L)<1$ the path is a monotone staircase inward, so convergence is one-sided; for $-1<f'(L)<0$ it spirals inward and the terms alternate about $L$; for $f'(L)>1$ the staircase runs outward and the sequence escapes; for $f'(L)<-1$ the spiral widens and the sequence oscillates without bound.

The structural consequence, which you should extract and use even without drawing anything: if $f$ is increasing the sequence is monotone, and if $f$ is decreasing the even and odd subsequences are each monotone, in opposite directions.

\begin{worked}
$x_1=1$, $x_{n+1}=\cos x_n$. Here $f=\cos$ is decreasing on $(0,1)$, so the cobweb spirals and the terms alternate around the crossing $L$ of $y=\cos x$ with $y=x$. Numerically $L=0.7390851332$, and $\cos'(L)=-\sin L=-0.6736$. The signed errors $x_n-L$ do alternate, with successive ratios $-0.762$, $-0.596$, $-0.716$, $-0.641$, $-0.693$, $-0.660$, settling slowly on $-0.6736$. Every consecutive pair of terms brackets the answer: $x_{2k}<L<x_{2k+1}$.
\end{worked}

\begin{trap}
Reading alternation as divergence. A spiral that tightens is convergence with a bonus: consecutive terms bracket the limit, which a monotone sequence never gives you. What signals divergence is a spiral that widens, that is $\abs{f'(L)}>1$.
\end{trap}

\begin{seealsob}Contraction and the derivative criterion; subsequence arguments; monotone convergence.\end{seealsob}

Two structural facts fall straight out of the picture and deserve to be stated without it. First, if $f$ is increasing the sequence is monotone, so the monotone convergence theorem is available and you need only a bound. Second, if $f$ is decreasing then $f\circ f$ is increasing, so the even-indexed and odd-indexed subsequences are each monotone, in opposite directions, and the question becomes whether their two limits agree. Those two sentences dispose of the qualitative behaviour of every one-step recursion with monotone $f$, which is most of them.

The cobweb is a picture of an inequality, and the inequality is the mean value theorem. Making it quantitative gives the single most useful rate statement in the subject.

\tech[Contraction and the derivative criterion]{Contraction and the $\abs{f'}<1$ criterion}{medium}
\cue A smooth recursion $x_{n+1}=f(x_n)$ with a known or computable fixed point $L$, and a question about the rate, or a monotonicity induction that refuses to close.
\method Subtract the fixed-point equation and apply the mean value theorem: $x_{n+1}-L=f(x_n)-f(L)=f'(\xi_n)(x_n-L)$ for some $\xi_n$ between $x_n$ and $L$. If $\abs{f'}\leq k<1$ on an interval $I$ containing $L$ with $f(I)\subseteq I$, then
\[\abs{x_n-L}\leq k^{\,n-1}\abs{x_1-L},\]
so convergence is geometric with ratio $\abs{f'(L)}$, and it happens for \emph{every} starting point in $I$. If $\abs{f'(L)}>1$ the fixed point repels and no starting point other than $L$ itself converges to it. If $f'(L)=0$ the linear term vanishes and convergence is faster than any geometric rate.

\begin{worked}
Two rates from the same idea.

\emph{Geometric.} For $x_{n+1}=\sqrt{2+x_n}$ with $L=2$, $f'(t)=\tfrac{1}{2\sqrt{2+t}}$ so $f'(2)=\tfrac14$. The errors from $x_1=1$ are $0.2679$, $0.06815$, $0.01711$, $0.004282$, $0.001071$, with successive ratios $0.2543$, $0.2511$, $0.2503$, $0.2501$. One digit per $1.66$ steps, forever.

\emph{Quadratic.} Newton's iteration for $\sqrt k$ is $x_{n+1}=\tfrac12\bigl(x_n+\tfrac{k}{x_n}\bigr)$, and $f'(t)=\tfrac12\bigl(1-\tfrac{k}{t^2}\bigr)$ vanishes at $t=\sqrt k$. With $k=2$ and $x_1=2$ the errors are $8.58\times10^{-2}$, $2.45\times10^{-3}$, $2.12\times10^{-6}$, $1.59\times10^{-12}$: each error is roughly the square of the previous one, which is what $f'(L)=0$ buys.

\emph{Repelling.} $x_{n+1}=x_n^2$ has fixed points $0$ and $1$, with $f'(0)=0$ and $f'(1)=2$. The first is superattracting and the second repels: from $x_1=0.99$ the sequence falls to $0$, from $x_1=1.01$ it runs to $\infty$, and only the exact value $1$ stays. No inequality argument will ever produce $1$ as a limit, and the fixed-point equation is powerless to say so.

Convert either statement into digits with logarithms. Geometric convergence with ratio $k$ buys $-\log_{10}k$ decimal digits per step, so the ratio $\tfrac14$ above gives $0.602$ digits per step, or one digit per $1.66$ steps, indefinitely. Quadratic convergence doubles the digit count each step, so the number of steps to reach any accuracy grows like $\log\log$ of it, which is why Newton's method is what calculators use.
\end{worked}

\begin{trap}
The borderline $\abs{f'(L)}=1$, where the criterion is silent and the sequence may converge, diverge, or crawl; the next technique handles it. A second failure is applying the estimate on an interval that is not forward invariant: $\abs{f'(L)}<1$ is a statement at a point, and without $f(I)\subseteq I$ nothing stops $x_2$ from landing outside $I$ where the bound does not hold.
\end{trap}

\begin{keynote}[THE DERIVATIVE IS THE RATE]
At an attracting fixed point the error is multiplied by $\abs{f'(L)}$ each step, so $\abs{f'(L)}$ is not a qualitative label but the actual convergence ratio. Its sign tells you the pattern (positive means one-sided, negative means alternating with brackets), its magnitude tells you the speed, and its vanishing tells you the speed is not geometric at all but quadratic.
\end{keynote}

\begin{seealsob}Sublinear decay when the derivative equals one; reading a recursion off the cobweb; monotone convergence.\end{seealsob}

The borderline case is not rare. It is exactly what happens whenever $f(x)=x-cx^p+\cdots$, which is the generic form of a recursion whose fixed point is at $0$ and whose first-order behaviour is the identity. Every such sequence converges, and none of them converges geometrically. The rate lives one order beyond the derivative, and there is a standard trick for extracting it.

\tech[Sublinear decay when the derivative equals one]{Sublinear decay: the rate when $f'(L)=1$}{hard}
\cue $x_{n+1}=x_n-c\,x_n^{p}+o(x_n^{p})$ with $c>0$ and $p>1$, after shifting the fixed point to $0$. The linearisation is the identity and gives no rate.
\method Do not study $x_n$; study the reciprocal power $y_n=x_n^{1-p}$, which grows linearly. Expanding,
\[x_{n+1}^{1-p}=x_n^{1-p}\bigl(1-cx_n^{p-1}+\cdots\bigr)^{1-p}=x_n^{1-p}+c(p-1)+o(1),\]
so the increments of $y_n$ tend to $c(p-1)$. By Stolz-Cesaro, $y_n/n\to c(p-1)$, hence
\[x_n\sim\bigl(c(p-1)\,n\bigr)^{-1/(p-1)}.\]
The decay is a power of $n$, never a power of a constant.

\begin{worked}
$x_1=1$, $x_{n+1}=\sin x_n$. Since $\sin t=t-\tfrac{t^3}{6}+O(t^5)$ we have $c=\tfrac16$ and $p=3$, so
\[x_n\sim\Bigl(\tfrac16\cdot2\cdot n\Bigr)^{-1/2}=\sqrt{\frac{3}{n}}.\]
The check: at $n=10^6$, $x_n\sqrt{n/3}=0.9999946$.

The same machinery on $x_{n+1}=x_n-x_n^2$ gives $c=1$, $p=2$, hence $x_n\sim1/n$, which you can also see directly from $1/x_{n+1}-1/x_n\to1$.

And on $x_{n+1}=\ln(1+x_n)$, from $\ln(1+t)=t-\tfrac{t^2}{2}+O(t^3)$ we read $c=\tfrac12$, $p=2$, so $x_n\sim\dfrac{2}{n}$. The check at $n=10^6$: $nx_n/2=1.0000023$. Three different recursions, three different constants, one method.
\end{worked}

\begin{trap}
Expanding only to first order, seeing $x_{n+1}\approx x_n$, and concluding either that the sequence fails to converge or that it converges geometrically. Both are wrong: it converges, and it converges like $n^{-1/(p-1)}$, which is far slower than any geometric rate. Halving the error in $x_{n+1}=\sin x_n$ costs a factor of four in $n$.
\end{trap}

\begin{seealsob}Contraction and the derivative criterion; Stolz-Cesaro; Cesaro means and the averaging principle.\end{seealsob}

\section{When the sequence does not converge, and when it is linear}

Everything so far has assumed the answer is yes. Some problems are posed precisely because the answer is no, and the standard demand then is a proof of non-existence, which is a different kind of object: you must produce evidence, not an estimate. An estimate can only ever say that the terms are somewhere; to deny a limit you must produce two places they keep returning to.

The right general framework is $\limsup$ and $\liminf$, which exist for every bounded sequence and are the largest and smallest subsequential limits. They convert the qualitative question into an equation: $\lim x_n$ exists exactly when $\liminf x_n=\limsup x_n$. In practice you rarely compute them; you exhibit two subsequences with different limits, which is the same thing done concretely.

\tech[Subsequence arguments and limsup liminf]{Subsequence arguments and $\limsup$, $\liminf$}{medium}
\cue You suspect a sequence does \emph{not} converge; or the recursion alternates, so the monotone argument is unavailable directly.
\method To disprove convergence, exhibit two subsequences with different limits; that is the entire proof, since every subsequence of a convergent sequence inherits the limit. To handle an alternating recursion, note that if $f$ is decreasing then $f\circ f$ is increasing, so the even terms $x_2,x_4,\ldots$ and the odd terms $x_1,x_3,\ldots$ are each monotone under $g=f\circ f$; find both limits and the original sequence converges if and only if they agree. Always available as a fallback: $\liminf x_n\leq\limsup x_n$, with equality exactly when the limit exists.

\begin{worked}
\emph{Disproof.} $a_n=\cos(n\pi/3)$ has $a_{6k}=1$ and $a_{6k+3}=-1$. Two subsequences, two limits, no limit.

\emph{Alternating recursion, convergent.} $x_1=0$, $x_{n+1}=\dfrac{6}{1+x_n}$. Here $f$ is decreasing, so the terms alternate about the positive fixed point: $L(1+L)=6$ gives $L=2$ or $L=-3$, and positivity selects $L=2$. Since $f'(2)=-\tfrac{6}{(1+2)^2}=-\tfrac23$ and $\abs{-\tfrac23}<1$, the even and odd subsequences both converge to $2$, so $x_n\to2$.

\emph{Alternating recursion, divergent.} $x_{n+1}=3-2x_n$ has the same fixed point $L=1$, but $f'=-2$, so from any $x_1\neq1$ the even and odd subsequences run to $+\infty$ and $-\infty$. The fixed-point equation is identical in both examples; only the derivative separates them.
\end{worked}

The $\limsup$ and $\liminf$ are worth keeping in reserve for sequences too irregular to describe by subsequences you can write down. For $a_n=\sin n$ the values are dense in $[-1,1]$, because the multiples of $1$ modulo $2\pi$ are equidistributed, so $\liminf a_n=-1$ and $\limsup a_n=1$; there is no single subsequence one would naturally exhibit, and the two-number summary is the entire available description.

\begin{trap}
Showing that one subsequence converges and calling the sequence convergent. The even terms of $(-1)^n$ converge to $1$ and prove nothing. Convergence of a single subsequence is evidence of nothing; convergence of the even and the odd subsequences to the \emph{same} value is a proof.
\end{trap}

\begin{seealsob}Monotone convergence; reading a recursion off the cobweb; contraction and the derivative criterion.\end{seealsob}

One family escapes all of this machinery, because it can be solved outright. When $f$ is linear and the recursion reaches back more than one step, the closed form is available and the limit question becomes a question about the size of two numbers.

\tech{Linear recurrences by characteristic root}{medium}
\cue $a_{n+1}=pa_n+qa_{n-1}$ with constant $p,q$, or a sequence built explicitly from Fibonacci or Lucas numbers.
\method Solve the characteristic equation $t^2=pt+q$. For distinct roots $\lambda_1,\lambda_2$ the general solution is $a_n=A\lambda_1^n+B\lambda_2^n$, with $A,B$ fixed by two initial values; for a repeated root $\lambda$ it is $a_n=(A+Bn)\lambda^n$. If $\abs{\lambda_1}>\abs{\lambda_2}$ the dominant term wins, so $a_{n+1}/a_n\to\lambda_1$ and $a_n^{1/n}\to\abs{\lambda_1}$.

\begin{worked}
The Fibonacci recursion $F_{n+1}=F_n+F_{n-1}$ has $t^2=t+1$, roots $\varphi=\tfrac{1+\sqrt5}{2}$ and $\bar\varphi=\tfrac{1-\sqrt5}{2}$, with $\abs{\bar\varphi}=0.618<\varphi$. Hence $F_n/F_{n-1}\to\varphi$. Then
\[\lim_{n\to\infty}\left(\frac{F_n}{F_{n-1}}-\frac{F_{n-1}}{F_n}\right)=\varphi-\frac1\varphi=\varphi-(\varphi-1)=1,\]
using the defining relation $\varphi^2=\varphi+1$ in the form $1/\varphi=\varphi-1$. At $n=60$ the double-precision value of the bracket is $1.0$ to machine accuracy.

Repeated root: $a_{n+1}=4a_n-4a_{n-1}$ has $t^2-4t+4=0$, so $\lambda=2$ twice and $a_n=(A+Bn)2^n$. With $a_0=0$, $a_1=2$ this is $a_n=n2^n$, and indeed $a_2=4\cdot2-4\cdot0=8=2\cdot2^2$.
\end{worked}

The dominant-root statement is worth unpacking, since it is the only place in the chapter where a limit is obtained by a cancellation rather than by an estimate. Writing $a_n=A\lambda_1^n+B\lambda_2^n$ with $\abs{\lambda_1}>\abs{\lambda_2}$ and $A\neq0$,
\[\frac{a_{n+1}}{a_n}=\lambda_1\cdot\frac{1+(B/A)(\lambda_2/\lambda_1)^{n+1}}{1+(B/A)(\lambda_2/\lambda_1)^{n}}\longrightarrow\lambda_1,\]
because $(\lambda_2/\lambda_1)^n\to0$. The convergence is geometric with ratio $\abs{\lambda_2/\lambda_1}$, which for Fibonacci is $0.618/1.618=0.382$, about $0.42$ decimal digits per step. The proviso $A\neq0$ matters: the Lucas-type sequence built to have $A=0$ is a pure multiple of $\lambda_2^n$ and its ratio tends to $\lambda_2$ instead.

One extension is worth knowing because it appears constantly in disguise. If the recursion carries a constant, $a_{n+1}=pa_n+q$ with $p\neq1$, subtract the fixed point $c=q/(1-p)$: the shifted sequence $a_n-c$ satisfies the homogeneous recursion $b_{n+1}=pb_n$, so $a_n=c+(a_1-c)p^{\,n-1}$. This converges exactly when $\abs{p}<1$, and to $c$, which is the contraction criterion for a linear $f$ with $f'=p$. From $a_1=1$, $a_{n+1}=2a_n+1$ we get $c=-1$ and $a_n=2^n-1$, a divergent case where the fixed-point equation nonetheless has the perfectly respectable root $-1$.

\begin{trap}
Writing $A\lambda^n+B\lambda^n$ for a repeated root. That is a one-parameter family wearing two letters, and it cannot match two independent initial conditions. The correct second solution carries the factor $n$.
\end{trap}

\begin{seealsob}Solving the fixed-point equation and selecting the admissible root; root and ratio limits for sequences; contraction and the derivative criterion.\end{seealsob}

\subsection{One recursion, all four questions}

It is worth seeing the whole apparatus applied once, in order, to a single problem, because in an examination all four questions are usually asked at once and the marks are allocated to the steps rather than to the answer.

\emph{Problem.} $x_1=0$, $x_{n+1}=\sqrt{6+x_n}$. Show the sequence converges, find the limit, give the rate, and say whether the terms over- or under-estimate the limit.

\emph{Existence.} $f(t)=\sqrt{6+t}$ is increasing on $[0,\infty)$. Bound: if $0\leq x_n<3$ then $x_{n+1}=\sqrt{6+x_n}<\sqrt9=3$, and $x_1=0<3$, so $x_n<3$ for all $n$. Monotone: $x_2=\sqrt6=2.449>0=x_1$, and $f$ increasing propagates this, so the sequence increases. Bounded and increasing, hence convergent.

\emph{Value.} Passing to the limit in the recursion, $L=\sqrt{6+L}$, so $L^2-L-6=(L-3)(L+2)=0$. Since $L\geq x_1=0$ we take $L=3$.

\emph{Rate.} $f'(t)=\dfrac{1}{2\sqrt{6+t}}$, so $f'(3)=\tfrac16$. The errors are $0.5505$, $0.09320$, $0.015574$, $0.0025967$, $0.00043282$, with successive ratios $0.1835$, $0.1693$, $0.1671$, $0.16674$, $0.166679$, settling on $\tfrac16$ as promised. That is $\log_{10}6=0.778$ decimal digits per step.

\emph{Pattern.} $f'(3)>0$, so the approach is one-sided, from below. Every term is therefore a rigorous lower bound for $3$, and no term is ever an upper bound. Had $f'(3)$ been negative, consecutive terms would have bracketed the limit instead, which is strictly more useful information.

Four questions, four different tools, and the only one that touches the fixed-point equation is the second. That proportion is the point of the chapter.

\section{Choosing the convergence argument}

The five recursion techniques above are not interchangeable, and the choice is forced by the shape of $f$, not by taste. The tree below is the order in which to test.

\begin{center}
\begin{tikzpicture}[node distance=5mm and 7mm]
\node[dtest] (a) {Linear with constant\\ coefficients?};
\node[dleaf, below left=9mm and 3mm of a] (b) {Characteristic roots;\\ dominant $\lambda$ wins};
\node[dtest, below right=9mm and 1mm of a] (c) {Is $f$ increasing on\\ the range?};
\node[dleaf, below left=9mm and 3mm of c] (d) {Monotone $+$ bounded,\\ then solve $L=f(L)$};
\node[dtest, below right=9mm and 1mm of c] (e) {Is $\abs{f'(L)}<1$ at the\\ candidate fixed point?};
\node[dleaf, below left=9mm and 3mm of e] (f) {Contraction: error\\ shrinks by $\abs{f'(L)}$};
\node[dnode, below right=9mm and 1mm of e] (g) {If $f'(L)=1$: expand one\\ order (sublinear decay).\\ Else: split even/odd.};
\draw[dedge] (a) -- node[dlab,left]{yes} (b);
\draw[dedge] (a) -- node[dlab,right]{no} (c);
\draw[dedge] (c) -- node[dlab,left]{yes} (d);
\draw[dedge] (c) -- node[dlab,right]{no} (e);
\draw[dedge] (e) -- node[dlab,left]{yes} (f);
\draw[dedge] (e) -- node[dlab,right]{no} (g);
\end{tikzpicture}
\end{center}

The tree is ordered by cost, not by power. Testing whether the recursion is linear costs one glance; testing whether $f$ is increasing costs one derivative; computing $f'(L)$ costs finding $L$ first, which is why it sits third even though it is the sharpest of the three. If two branches apply, take the earlier one: a linear recursion should never be attacked with a contraction estimate, and an increasing $f$ with an obvious bound should never be attacked with the mean value theorem.

A caution about the third question. Writing down $\abs{f'(L)}$ presumes you already know $L$, which presumes the fixed-point equation has been solved and the right root chosen. That is legitimate here in a way it was not in the opening section, because the derivative test is a genuine existence proof: if $\abs{f'}\leq k<1$ on a forward-invariant interval containing $L$, convergence follows for every start in that interval, and $L$ is then the limit rather than a candidate for it. The order of operations is still solve, then verify, then conclude, but here the verification does the work that the induction did before.

Two remarks on using it honestly. First, the leftmost yes branch at the second question still owes you a bound; ``$f$ is increasing'' delivers monotonicity for free but nothing else. Second, the tree tells you which existence argument to attempt, and the value of the limit is in every case obtained afterwards from $L=f(L)$. There is no branch on which the fixed-point equation is the argument.

\section{Discrete asymptotics}

The rest of the chapter changes the question. Instead of a recursion whose limit is a number, we are given an explicit but unwieldy $a_n$ (a partial sum, a factorial, an $n$th root) and asked for its size. The tools here are transfer principles: each one converts a limit you cannot compute into one you can, at the cost of a hypothesis that must be checked.

Before the tools, the hierarchy. Every asymptotic argument in this section is ultimately a statement about where a quantity sits in the chain
\[1\ll\ln\ln n\ll\ln n\ll n^{\alpha}\ll c^{\,n}\ll n!\ll n^{n}\qquad(\alpha>0,\;c>1),\]
where $u_n\ll v_n$ means $u_n/v_n\to0$. Each step is strict and each gap is enormous: $\ln n$ passes $100$ only past $n=10^{43}$, while $n!$ passes $10^{100}$ before $n=70$. When a limit has the form $u_n/v_n$ and the two sides sit at different links of the chain, the answer is $0$ or $\infty$ and no technique is needed. The techniques below are for the cases where they sit at the same link and the constant, or the polynomial correction, is what is being asked for.

One transfer that belongs here but is developed elsewhere in this book deserves a signpost, because it is the first thing to try on a limit of the shape $\tfrac1n\sum_{k\le n}g(k/n)$: that is a Riemann sum, and its limit is $\int_0^1g$. The signature is a sum whose summand depends on $k$ only through the ratio $k/n$. When the summand instead depends on $k$ and $n$ separately, the Riemann reading fails and the techniques below are what remain.

The first is the discrete L'Hopital, and it earns its keep on any ratio whose numerator is a sum. Differences are to sequences what derivatives are to functions, and a sum has a trivially simple difference.

\tech{Stolz-Cesaro}{medium}
\cue A ratio $a_n/b_n$ where $b_n$ is strictly monotone and unbounded: the discrete $\infty/\infty$. Almost always $a_n$ is a partial sum, or something whose consecutive difference is much simpler than itself.
\method If $b_n$ is strictly increasing and unbounded and
\[\lim_{n\to\infty}\frac{a_{n+1}-a_n}{b_{n+1}-b_n}=L,\]
then $\lim a_n/b_n=L$ too. This is the exact sequence analogue of L'Hopital's rule and, like it, is a one-way implication: the difference ratio may fail to have a limit while the original ratio converges perfectly well.

\begin{worked}
\[\lim_{n\to\infty}\frac{1^p+2^p+\cdots+n^p}{n^{p+1}}\qquad(p>0).\]
Take $a_n$ the sum and $b_n=n^{p+1}$, which is strictly increasing and unbounded. Then
\[\frac{a_{n+1}-a_n}{b_{n+1}-b_n}=\frac{(n+1)^p}{(n+1)^{p+1}-n^{p+1}}.\]
The denominator is $(p+1)n^p+O(n^{p-1})$ by the binomial theorem, so the ratio tends to $\dfrac{1}{p+1}$, and that is the answer. At $p=3$ and $n=10^4$ the partial value is $0.2500500$ against $\tfrac14$.

A second, more instructive instance: $a_n=H_n=\sum_{k\le n}1/k$ and $b_n=\ln n$. The difference ratio is $\dfrac{1/(n+1)}{\ln(1+1/n)}\to1$, so $H_n/\ln n\to1$. Note how slowly: at $n=10^7$ the ratio is still $1.0358$, because $H_n=\ln n+\gamma+o(1)$ and the correction only dies like $\gamma/\ln n$. Stolz-Cesaro gives the limit and says nothing about the rate.
\end{worked}

There is a second, quieter use of the same theorem. Because it turns a sum into its summand, it is the standard way to convert a statement about increments into a statement about size: if $a_{n+1}-a_n\to c$ then $a_n/n\to c$, so $a_n\sim cn$. That is the step used to extract the rate in the sublinear decay technique, applied there to $y_n=x_n^{1-p}$, and it is the reason that technique works at all.

One more instance, to show the range. Let $a_n=\ln(n!)$ and $b_n=n\ln n$. The difference ratio is
\[\frac{\ln(n+1)}{(n+1)\ln(n+1)-n\ln n},\]
whose denominator is $\ln n+1+o(1)$, so the ratio tends to $1$ and $\ln(n!)\sim n\ln n$. That is the crude half of Stirling's approximation, obtained with no analysis at all beyond a difference and a limit. What Stolz-Cesaro cannot give you is the next term, the $-n$; for that you need the real thing.

\begin{trap}
Applying it when $b_n$ is not monotone; the hypothesis is not decoration. And, exactly as with L'Hopital, concluding that the original ratio has no limit because the difference ratio has none. Take $a_n=(-1)^n$, $b_n=n$: the difference ratio oscillates between $\pm2$ and has no limit, while $a_n/b_n\to0$.
\end{trap}

\begin{seealsob}Cesaro means and the averaging principle; sublinear decay when the derivative equals one; root and ratio limits for sequences.\end{seealsob}

It is worth being explicit about the direction of the implication, because it is the same asymmetry as L'Hopital and it is misused the same way. Stolz-Cesaro says: if the difference ratio converges, so does the original. It never says the reverse. In practice this means Stolz-Cesaro can only ever produce an answer; it can never justify the sentence ``the limit does not exist''. For that you need the subsequence machinery.

The special case $b_n=n$ deserves its own name, because it says something structural: averaging is a smoothing operation. It preserves limits and it can manufacture them out of oscillation, which is why the implication runs one way only.

\tech{Cesaro means and the averaging principle}{medium}
\cue An arithmetic mean $\tfrac1n\sum_{k\le n}a_k$, or a geometric mean $(a_1a_2\cdots a_n)^{1/n}$, or any expression that is secretly one of these.
\method If $a_n\to L$ then the arithmetic means converge to $L$; if in addition $a_n>0$ then the geometric means converge to $L$ as well. (The geometric statement is the arithmetic one applied to $\ln a_n$.) The averaging is a contraction on oscillation: it can only destroy irregularity, never create it.

\begin{worked}
\[\lim_{n\to\infty}\frac{n^2}{H_n\displaystyle\sum_{k=1}^{n}\frac{k}{H_k}}.\]
Let $S_n=\sum_{k\le n}k/H_k$. Apply Stolz-Cesaro with $b_n=n^2/H_n$, which is strictly increasing and unbounded. Since $H_{n+1}-H_n=1/(n+1)$ is negligible against $H_n$, the increment of $b_n$ is $\bigl(2n+1\bigr)/H_{n+1}+o(n/H_n)$, while the increment of $S_n$ is $(n+1)/H_{n+1}$. The ratio tends to $\tfrac12$, so $S_n\sim\dfrac{n^2}{2H_n}$ and the required limit is
\[\frac{n^2}{H_n\cdot n^2/(2H_n)}=2.\]
The numerical trend is $1.891$, $1.913$, $1.928$, $1.938$ at $n=10^4,10^5,10^6,10^7$: monotone increasing toward $2$ with an $O(1/\ln n)$ correction, which is precisely what an $H_n$ in the denominator predicts.
\end{worked}

The geometric version is the one that does real work, because so many discrete limits are products in disguise. A second instance: to evaluate
\[\lim_{n\to\infty}\left(\frac{(2n)!}{n!\;n^{n}}\right)^{1/n},\]
set $a_n=(2n)!/(n!\,n^n)$ and compute the ratio, $\dfrac{a_{n+1}}{a_n}=\dfrac{2(2n+1)}{n+1}\Bigl(\dfrac{n}{n+1}\Bigr)^{n}\to\dfrac{4}{\eu}$. The geometric-mean principle then delivers the root limit $4/\eu=1.47152$, confirmed numerically as $1.47152$ at $n=2\times10^{5}$. Doing this directly, by expanding both factorials, is a much longer computation for the same answer.

\begin{trap}
Reading the converse. If the Cesaro means converge, the sequence need not: $a_n=(-1)^n$ has means tending to $0$ and no limit at all. Averaging is where information goes to die, so you may push limits through it but never pull them back.
\end{trap}

\begin{seealsob}Stolz-Cesaro; root and ratio limits for sequences; Stirling's approximation.\end{seealsob}

The geometric-mean half of the last technique has a consequence that is used far more often than the technique itself: it converts a ratio limit into a root limit. Since ratios of consecutive terms are usually easy and $n$th roots are usually not, the transfer runs in the useful direction.

\tech{Root and ratio limits for sequences}{medium}
\cue An $n$th root of a rapidly growing quantity (a factorial, a product, a power), or a ratio of consecutive terms that simplifies.
\method If $a_n>0$ and $\lim a_{n+1}/a_n=L$, then $\lim\sqrt[n]{a_n}=L$. The proof is the geometric-mean principle applied to the telescoped product $a_n=a_1\prod_{k<n}(a_{k+1}/a_k)$. The converse is false, so the root limit is strictly the stronger statement: it exists in cases where the ratio limit does not.

\begin{worked}
\[\lim_{n\to\infty}\frac{n}{\sqrt[n]{n!}}.\]
Set $a_n=n^n/n!$, so that $n/\sqrt[n]{n!}=\sqrt[n]{a_n}$. The ratio is
\[\frac{a_{n+1}}{a_n}=\frac{(n+1)^{n+1}}{(n+1)!}\cdot\frac{n!}{n^n}=\frac{(n+1)^n}{n^n}=\Bigl(1+\frac1n\Bigr)^{n}\to\eu.\]
Hence the root limit is $\eu$ as well. At $n=2\times10^5$ the value is $2.71819$ against $\eu=2.718282$.
\end{worked}

The transfer also explains a fact used constantly without comment: $\sqrt[n]{n}\to1$, $\sqrt[n]{n!}\to\infty$ like $n/\eu$, and $\sqrt[n]{c}\to1$ for any fixed $c>0$. Each is a ratio computation of one line. And it is the sequence statement underneath the root and ratio tests for series, where exactly the same asymmetry appears: the root test decides every series the ratio test decides, and some that it does not.

A second instance, worth doing because the direct approach is genuinely unpleasant: $\lim\sqrt[n]{\binom{2n}{n}}$. The ratio is
\[\frac{\binom{2n+2}{n+1}}{\binom{2n}{n}}=\frac{(2n+1)(2n+2)}{(n+1)^2}=\frac{2(2n+1)}{n+1}\to4,\]
so the root limit is $4$. Numerically $\sqrt[n]{\binom{2n}{n}}=3.998$ at $n=10^4$, approaching $4$ from below at the leisurely rate $4\,(\pi n)^{-1/(2n)}$, which is what the next technique will make precise.

\begin{trap}
Concluding that the root limit fails because the ratio limit does. Take $a_n=2^{\,n+(-1)^n}$: the ratio $a_{n+1}/a_n=2^{\,1-2(-1)^n}$ alternates between $8$ and $\tfrac12$ and has no limit, yet $\sqrt[n]{a_n}=2^{1+(-1)^n/n}\to2$ without difficulty. The ratio test is a sufficient condition, not a diagnostic.
\end{trap}

\begin{seealsob}Cesaro means and the averaging principle; Stirling's approximation; linear recurrences by characteristic root.\end{seealsob}

The previous technique gets the exponential scale of $n!$ right and throws away everything else. When the normalisation is finer, you need the full asymptotic, including the constant, and that is Stirling.

\tech{Stirling's approximation}{medium}
\cue A factorial under a root, inside a logarithm, or divided by a power of $n$. Also central binomial coefficients, which are factorials in disguise.
\method Use
\[n!\sim\sqrt{2\pi n}\left(\frac{n}{\eu}\right)^{n},\qquad\text{equivalently}\qquad \ln(n!)=n\ln n-n+\tfrac12\ln(2\pi n)+o(1).\]
Take logarithms \emph{first}, substitute, keep only the terms that survive the normalisation you are given, then exponentiate. Deciding which terms survive is the entire skill.

\begin{worked}
Three normalisations of the same asymptotic.

\emph{The constant dies.} $\dfrac{\sqrt[n]{n!}}{n}$: taking logs, $\tfrac1n\ln(n!)-\ln n=-1+\tfrac{1}{2n}\ln(2\pi n)+o(1/n)\to-1$, so the limit is $\eu^{-1}=0.36788$. Verified: $0.36789$ at $n=2\times10^5$. The factor $\sqrt{2\pi n}$ contributes $(2\pi n)^{1/(2n)}\to1$ and vanishes.

\emph{The constant is the whole answer.} $\dfrac{n!\,\eu^{n}}{n^{n}\sqrt n}\to\sqrt{2\pi}=2.50663$. Here the exponential parts cancel exactly by construction and only the constant is left. At $n=1000$ the value is $2.50684$.

\emph{A mixed case.} $\dbinom{2n}{n}=\dfrac{(2n)!}{(n!)^2}\sim\dfrac{\sqrt{4\pi n}\,(2n/\eu)^{2n}}{2\pi n\,(n/\eu)^{2n}}=\dfrac{4^{n}}{\sqrt{\pi n}}$: the constant survives but in the combination $\sqrt{\pi}$, not $\sqrt{2\pi}$. At $n=1000$ the ratio of the two sides is $0.999875$.
\end{worked}

A fourth normalisation, which appears whenever double factorials do: $(2n)!!=2^nn!$ and $(2n-1)!!=\dfrac{(2n)!}{2^nn!}$, so
\[\frac{(2n-1)!!}{(2n)!!}=\frac{(2n)!}{4^n(n!)^2}=\frac{1}{4^n}\binom{2n}{n}\sim\frac{1}{\sqrt{\pi n}}.\]
At $n=10^4$ the ratio of the two sides is $0.9999875$. This single asymptotic decides an entire family of series, and it is the reason those series sit exactly at the ratio test's blind spot.

The log form is also the practical form. The number of decimal digits of $n!$ is $\floorb{\log_{10}(n!)}+1$, and Stirling gives $\log_{10}(100!)\approx(100\ln100-100+\tfrac12\ln(200\pi))/\ln10=157.9696$ against the exact $157.9700$, so $100!$ has $158$ digits. Four correct significant figures from a formula with one term of error.

If more accuracy is wanted, the next term of the asymptotic series is available and is worth memorising because it is so cheap:
\[n!=\sqrt{2\pi n}\left(\frac{n}{\eu}\right)^{n}\left(1+\frac{1}{12n}+O(n^{-2})\right).\]
At $n=10$ the uncorrected formula gives $3598696$ against the exact $3628800$, an error of $0.83$ per cent; the single correction factor $1+\tfrac{1}{120}$ brings it to $3628685$, an error of $0.003$ per cent. Even at $n=5$ the corrected value is $119.99$ against $120$. Stirling is not merely asymptotic; it is accurate immediately.

\begin{trap}
Keeping $\sqrt{2\pi n}$ when the normalisation kills it, or dropping it when it survives. The rule is mechanical once you take logarithms: the factor contributes $\tfrac12\ln(2\pi n)$, and you ask whether that term is $o$ of the normalisation. In $\sqrt[n]{n!}/n$ it is divided by $n$ and dies; in $n!\eu^n/(n^n\sqrt n)$ it is all that remains.
\end{trap}

\begin{seealsob}Root and ratio limits for sequences; Cesaro means and the averaging principle; Stolz-Cesaro.\end{seealsob}

\begin{keynote}[SUMS OR PRODUCTS]
Stolz-Cesaro and the Cesaro means handle sums, because differencing a sum is trivial. Root and ratio limits handle products, because a logarithm turns a product into a sum. Stirling handles the one product that occurs more often than all the others together. When a discrete limit resists, the productive question is not ``which clever manipulation'' but ``is the object underneath a sum or a product'', because that single distinction selects the transfer principle.
\end{keynote}

One last connection back to the first half of the chapter. The asymptotic techniques and the recursion techniques are not two subjects. The sublinear decay technique is Stolz-Cesaro applied to a reciprocal power; the contraction estimate is the statement that a certain sequence of errors is geometric, which is a statement about where it sits in the growth hierarchy; the linear recurrence technique is the observation that a sum of two exponentials is asymptotically the larger one. In every case the recursion supplies increments and the asymptotic machinery converts increments into size. If you find yourself with a recursion you cannot solve and a rate you need, difference something and apply Stolz-Cesaro; it is the move that works more often than any other in this chapter.

Read the matrix by shape, not by topic. The left column lists what is visible on the page before any work has been done: the form of $f$, the form of $a_n$, the presence of a factorial or a partial sum. The middle column is the technique that shape selects, and the right column is what to do when the hypothesis of that technique fails to check out, which is usually the same technique's nearest neighbour rather than a fresh idea.

\section{Recognition matrix}
\begin{recog}{@{}p{0.30\textwidth}p{0.34\textwidth}p{0.28\textwidth}@{}}
\rowcolor{mist}\mh{If you see} & \mh{Reach for} & \mh{If that fails} \\
\midrule
\endhead
$x_{n+1}=f(x_n)$, $f$ increasing & Monotone $+$ bounded induction, then $L=f(L)$ & Contraction estimate on an invariant interval \\
$x_{n+1}=f(x_n)$, $f$ decreasing & Even and odd subsequences under $f\circ f$ & Cobweb, then $\abs{f'(L)}$ \\
Nested radical, continued fraction, tower & $L=f(L)$, then discard roots outside the range & Monotone convergence for the bound \\
Two roots both satisfying $L=f(L)$ & The proved bound on $x_n$ selects & Sign and monotonicity of the sequence \\
$x_{n+1}=x_n-cx_n^{p}$, $p>1$ & Reciprocal power: $x_n^{1-p}\sim c(p-1)n$ & Stolz-Cesaro on $x_n^{1-p}$ \\
$a_{n+1}=pa_n+qa_{n-1}$ & Characteristic roots; dominant $\lambda$ & Repeated root gives $(A+Bn)\lambda^n$ \\
Partial sum over a power of $n$ & Stolz-Cesaro & Recognise it as a Riemann sum \\
$\tfrac1n\sum a_k$ or $(\prod a_k)^{1/n}$ & Cesaro means & Stolz-Cesaro directly \\
$\sqrt[n]{a_n}$ with $a_n$ a product & Ratio-to-root transfer & Logarithms plus Cesaro means \\
$n!$ under a root, log, or power & Stirling, in log form & Ratio-to-root transfer if the constant dies \\
Suspected divergence & Two subsequences, two limits & $\liminf<\limsup$ \\
\end{recog}
